\documentclass[../main.tex]{subfiles}
\begin{document}
%==========================================TIÊU ĐỀ TẬP========================================
\begin{table}[h]
    \begin{adjustwidth}{-1cm}{}
        \begin{tabular}{>{\centering\arraybackslash}p{6cm}|>{\raggedright\arraybackslash}p{12.5cm}}
            \multirow{5}{6cm}
            % Cột bên trái: ÉPISODE và số tập
            {\\[-40pt]
                \flaregothic\fontsize{20pt}{20pt}\selectfont \centering
                \color{eptype}ÉPISODE\color{black}\\
                \burbank\fontsize{72pt}{72pt}\selectfont
                \color{epnum}32\color{black}\\[6pt]
                \cafeteria\fontsize{20pt}{20pt}\selectfont\color{epnum}(3-3)\color{black}
                \\[18pt]
                \flaregothic\fontsize{18pt}{18pt}\selectfont
                \color{epattr}ÉPISODE PERDU
                \color{black}
            }
             &
            % Dòng thứ nhất: tên tập tiếng Nhật
            {\fotskip\fontsize{18pt}{18pt}\selectfont \ruby{ア}{あ}\ruby{イ}{い}\ruby{ド}{ど}\ruby{ル}{る}の\ruby{手}{て}\ruby{料}{りょう}\ruby{理}{り}といえば……【\ruby{死}{し}\ruby{者}{しゃ}\ruby{蘇}{そ}\ruby{生}{せい}】
            }                                                                                      \\[6pt]
            % Dòng thứ hai: tên tập tiếng Anh
             & {\cafeteria\fontsize{18pt}{18pt}\selectfont When Cooking Gets Fishy (Gets Reborned!)}                                                                             \\[4pt]
            % Dòng thứ ba: tên tập tiếng Việt
             & {\vnmsans\fontsize{18pt}{18pt}\selectfont Bát cơm bò và người cá giao hàng}                                                                                 \\[8pt]
            % Dòng thứ tư: Nhân vật xuất hiện
             & {\caviar{\fontsize{14pt}{14pt}\selectfont Nhân vật: \emp{kuon1} \emp{noel} \emp{choco} \emp{haato}}
               }
            \\[8pt]
            % Dòng cuối cùng: Thời gian diễn ra của tập (thường sẽ là ngày phát hành)
             & {\continuum\fontsize{16pt}{16pt}\selectfont Ngày phát hành: 15/12/2019}                                                                           \\[18pt]
        \end{tabular}
    \end{adjustwidth}
\end{table}
%=========================================TRUYỆN CHÍNH========================================
\color{black}

\txtbx[2]{{\color{vol} \uline{Tóm tắt:}} Cơn đói cồn cào bao trùm văn phòng cho đến khi Haato mang về một con cá ngừ khổng lồ từ trò chơi xổ số. Tuy nhiên, con cá ngừ này lại có khả năng ``tiến hóa'' kỳ dị, dẫn đến một cuộc rượt đuổi hỗn loạn và một bát cơm bò (Gyudon) có một không hai được giao bởi chính\dots\ con cá đó.}

Một buổi trưa giữa tháng Mười Hai. Tuyết nhẹ rơi lác đác ngoài cửa sổ, vẽ nên một khung cảnh mùa đông yên bình. Thế nhưng, bên trong văn phòng \hll, sự yên bình đó đang bị đe dọa bởi những cái bụng đói cồn cào.

\img{3-3a}

Noel không ngừng kêu ca, tay xoa bụng đầy vẻ thèm thuồng: ``Mình muốn ăn món gyudon\footnote{Món cơm ăn với thịt bò của Nhật Bản.} quá\dots\ chỉ nghĩ đến lớp thịt bò mềm mại thôi mà đã chảy nước miếng rồi!''

Choco ngồi trên ghế, ánh mắt lộ vẻ bồn chồn không kém. Cô khẽ thở dài, tay ôm chặt con khỉ bông như để phân tán cơn đói: ``Bình tĩnh lại nào Noel-sama. Cô cũng chưa có gì bỏ vào bụng từ sáng đây này.''

Riêng Kuon vẫn bình thản ngồi làm việc, nét mặt điềm nhiên như không có gì xảy ra. Có lẽ ``tiểu sử nhịn ăn sáng'' gia truyền đã giúp cậu có sức chịu đựng bền bỉ hơn hẳn. Tuy nhiên, thi thoảng cậu vẫn lén liếc nhìn hai cô gái với vẻ ái ngại: ``\textit{Nhịn đói giữa cái tiết trời này đúng là cực hình thật.}''

Bất chợt, cánh cửa văn phòng bật mở. Haato lao vào với vẻ mặt hớn hở tột độ, tay ôm một vật thể khổng lồ: ``Chị trúng được con cá ngừ ở trò xổ số này!!''

\img{3-3b}

Cả ba người trong phòng đồng loạt đứng hình. Trước mặt họ là một con cá ngừ to lớn, dài gần bằng thân người của Haato. 

``Nguyên\dots\ con cá này luôn á!?'' Kuon trố mắt kinh ngạc.

Choco cũng không tin vào mắt mình: ``Họ cứ thế đưa cho em nguyên con như vậy thôi sao? Không làm sạch hay gì luôn?''

Cậu Tổng nhíu mày, nhìn con cá ``còn nguyên zin'' từ đầu đến vây với vẻ nghi hoặc cực độ. Nhưng Haachama chỉ cười toe toét, khẳng định chắc nịch rằng: ``Đúng thế! Người ta cứ đưa cho em nguyên con như vậy thôi á!'' 

Cô lập tức đặt con cá lên bàn, tuyên bố: ``Mà, chắc mọi người cũng đói rồi nhỉ? Để em nấu cái gì đó nhé?''

Vừa nghe đến việc Haachama nấu ăn, sắc mặt cậu Tổng lập tức biến đổi. Cậu vội vàng ngăn cản: ``Ờ, thôi, để tí nữa anh vào anh nấu cho. Em cứ để con cá đó ở đấy đã-''

Nhưng vô ích. Với sự bướng bỉnh đã trở thành thương hiệu, Haachama hí hửng ôm con cá vào bếp, bỏ lại ba con người đang nhìn theo với ánh mắt đầy lo âu và sợ hãi. ``Anh cứ ngồi đó đi, để em nấu cho! Không phải phờ-lo!''

Cậu Tổng tựa lưng vào ghế, thở dài bất lực: ``Hy vọng là không phải như cái hồi em ấy làm bánh mì kẹp tháng trước\dots\''

Choco ngơ ngác, tò mò hỏi cậu: ``Ý anh là sao, Kuon-sama?'' Cậu nhắm mắt, như muốn quên đi một ký ức kinh hoàng, rồi chậm rãi kể: ``Hôm đó chỉ có mình anh trong văn phòng. Tự dưng không biết vì lý do gì, Haato-chan nổi hứng muốn nấu ăn\dots''

Nàng Hiệp sĩ chống cằm, mắt cô sáng lên: ``Rồi sao nữa ạ?''

Kuon nhún vai, khuôn mặt đầy bất lực: ``Kết quả là không ăn được. Bánh mì thì ăn như than, thịt thì cứng như đá cắn gãy cả răng, rau thì còn nguyên mùi đất\dots\ Mang về cho chó nó còn không thèm nữa là.''

Nghe đến đây, cả Noel và Choco đều sợ xanh mặt. Choco lẩm bẩm, ôm khỉ bông chặt hơn: ``C-Chắc không đến mức đó đâu nhỉ?'' Noel thì cười gượng, nhưng mắt vẫn dán về phía bếp, nơi mà Haachama đang ``tung hoành'': ``Haachama-chan thì khó mà đoán được\dots\ Nhưng mà nếu là gyudon thì chắc không sao đâu, nhỉ?''

Cậu nhìn nàng Hiệp sĩ một cách đầy hoài nghi, lẩm bẩm đủ để hai cô gái nghe thấy: ``Tốt nhất là nên chuẩn bị tinh thần ăn\dots\ sushi sống từ con cá này thì hơn.''

Tiếng hét thất thanh từ trong bếp đột ngột vang lên, cắt ngang câu chuyện kinh dị của Kuon. Noel, do cơn đói đã lên đến đỉnh điểm, liền bật dậy hét lớn: ``\textbf{Gyudon!?}''

Hai người kia lập tức ngó vào bếp, chỉ để chứng kiến một cảnh tượng không tưởng: Con cá ngừ mà Haachama định làm thịt đang cựa quậy dữ dội trên sàn.

\img{3-3c1}

``Làm thế nào mà nó vẫn sống thế này!?'' cậu thảng thốt.

Nàng Nữ sinh trung học, đang loay hoay nhặt lại con dao bếp rơi dưới đất, lúng túng giải thích: ``Em định nấu nó bằng bếp từ cảm ứng thôi mà! Hay là nó bị điện giật chăng?''

``Thử kiểm tra xem bếp có bị rò điện không? Nếu không thì-'' Nhưng chưa kịp để cậu nói hết câu, con cá ngừ đột ngột vùng lên, mọc ra bốn chi giống hệt con người và bắt đầu chạy loạn xạ.

\img{3-3c2}

``\textbf{Cái\dots!? Con cá này là người à!?}'' 

Ngay sau đó, con cá bất ngờ cắm đầu chạy khỏi bếp, để lại cả ba con người ngẩn ngơ không nói nổi câu nào.

``Này! Đứng lại!!'' nàng Song tính định đuổi theo nó, nhưng nó chạy quá nhanh khi cô còn chưa kịp đứng dậy.

Nó lao qua trước mặt Choco và Noel, một người vẫn đang ôm chặt con khỉ bông, còn người kia thì gần như lả đi vì đói. Thế rồi, con cá ngừ phá tung cửa sổ rồi biến mất trong màn tuyết trắng ngoài trời.

\img{3-3d}

Mặt Kuon tái mét: ``Thế là bay thêm một cái cửa sổ nữa\dots\ A-san mà biết thì chúng ta xong đời.''

Haachama thất thểu bước ra khỏi bếp, mặt đầy thất vọng. Cô áp sát người vào tường, giọng như đang oán trách vận mệnh: ``Không\dots\ Cá ngừ của mình\dots''

Kuon thở phào, mỉm cười nhẹ nhõm: ``Ít nhất thì chúng ta không phải ăn món em ấy làm rồi\dots'' Trong khi đó, Noel, người ban nãy còn lả đi vì đói, bỗng chốc tràn đầy năng lượng. Cô cầm điện thoại, cười tươi rói và đưa lên khoe: ``Thôi thì, mình đã đặt gyudon rồi!''

\ct

Một lát sau, tiếng chuông cửa vang lên. Noel lập tức dùng cây chùy sắt của mình phá tung cánh cửa để đón nhân viên giao hàng.

``Có cần thiết phải bạo lực thế không, Hiệp sĩ ơi? Hay là em đói quá nên mất trí rồi?'' Kuon chống cằm, nhíu mày nhìn cô nàng với hành động không giống người thường kia.

\img{3-3e}

Nhưng người (hay vật) đứng ở cửa không phải là shipper bình thường. Đó chính là con cá ngừ ban nãy, hiện đang bưng một bát gyudon khổng lồ trên tay.

\img{3-3f}

Nhìn bát cơm bò trên tay, nàng Hiệp sĩ vô thức hét lớn: ``\textbf{Gyudon!?}'' trong khi nàng Y tá và cậu Tổng thì ngơ ngác nhìn kẻ đang đứng ở trước cửa kia.

``Chuyện quái gì đang xảy ra thế\dots?'' Kuon nhăn mặt. Theo sau, Choco cũng tiếp lời: ``Hay có lẽ vì chúng ta đói quá chăng nên nhìn gà hóa cuốc\dots?''

Không để ai kịp phản ứng, con cá thu mình lại, tự đặt mình lên bát cơm, biến thành một món ``cơm bò độn cá ngừ tươi sống'' có một không hai.

Cả văn phòng chìm vào im lặng. Noel, trong cơn đói đỉnh điểm, chỉ biết thốt lên đầy phấn khích: ``\textbf{Kaidendon!!!}'' (Cơm bò hải sản!!!)

\img{3-3g}

Haato và Choco đứng nhìn bát cơm với vẻ mặt đầy hỗn độn, không biết nên khóc hay nên cười trước món ăn kỳ quái này. Còn Kuon nhìn chỉ chống hông bất lực nhìn cô nàng Hiệp sĩ đang phấn khích với bữa cơm trưa của mình.

%==========================================TRUYỆN PHỤ=========================================
\omk

Cuối cùng, Kuon vẫn phải tự mình xuống bếp để xử lý con cá ngừ (vốn đã thôi không chạy nhảy nữa). Cậu ngăn không cho Haato vào bếp, vừa lắc đầu vừa nói: ``Anh bảo rồi, để anh lo. Nếu để em nấu thì nó chạy thật đấy!'' Rõ ràng đó chỉ là một cái cớ, vì cậu không muốn cô nàng Nữ sinh kia động vào bất cứ thứ gì vào bếp cả.

Choco liền gật đầu lia lịa đồng ý: ``Kuon-sama nói đúng đó. Để Haato nấu thì chắc con cá này lại sống lại thật mất!''

Noel lanh lảnh: ``Đúng vậy! Kuon-senpai nấu thì chắc chắn sẽ ngon rồi!''

Haachama ngồi phịch xuống ghế, mặt phụng phịu: ``Em có phải vụng về đến mức đó đâu\dots\ Mà cá ngừ sống có khi ăn được mà, đúng không Choco-sensei?''

Nàng Y tá nhún vai, nhưng cũng không quên lườm cô, bày tỏ sự lo lắng: ``Cô không muốn thử đâu. Hôm nay là cá ngừ, lỡ ngày mai em nhặt về con cá sấu thì sao?''

Nghe vậy, Haato cười trừ, nhưng nàng Hiệp sĩ thì cười lớn: ``Haato-chan mà nấu cá sấu thì\dots\ chắc là dị lắm đây!''

\ct

Sau một hồi chế biến, Kuon bê ra bàn những đĩa cá ngừ áp chảo thơm phức cùng những tô gyudon nóng hổi. Mùi thơm ngào ngạt khiến cả ba cô gái không thể cưỡng lại. ``Xong rồi đây. Các em ăn nhanh còn làm việc.''

Mùi thơm từ đĩa cá khiến cả ba cô gái không thể cưỡng lại.  Noel vừa ăn vừa xuýt xoa: ``Đây đúng là thiên đường! Gyudon và hải sản, món ưa thích của mình! Lần sau có cá ngừ, mình chỉ mong Kuon-senpai nấu thôi!''

Cô ăn một miếng cá, khuôn mặt lộ rõ vẻ hạnh phúc đến mức Haachama phải nhíu mày: ``Ơ\dots\ Thích thế cơ à? Nhưng mà\dots\ sao em thấy hơi nhạt nhỉ? Hay tại em thích cá ngừ sống hơn?''

Choco, đang thong thả thưởng thức, quay sang nhìn cô và đáp: ``Thế thì lần sau em thử tự làm sashimi đi, Haato-sama, nhưng nhớ là phải dùng cá tươi thật nhé, không phải cá nhảy nhót khắp nơi đấy.''

Kuon bật cười: ``Sashimi? Cứ nghĩ đến cảnh Haachama-chan cắt cá mà anh đã muốn té trước rồi.'' Haato chu môi, phản bác: ``Ít nhất em còn biết cách xử lý cá hơn là anh tưởng!''

Cô nói như vậy, nhưng chẳng ai tin cô cả. Họ chỉ cần nhìn cách cậu Tổng phản ứng với các món ăn của cô là họ hiểu cô ấy nấu tệ đến cỡ nào.

\ct

Sau bữa ăn, Kuon dọn dẹp bát đĩa trong khi ba cô gái ngồi nghỉ. Choco ngả người trên ghế, tay ôm con khỉ bông, nói bâng quơ: ``Phải công nhận, Kuon-sama nấu ăn không tệ.''

Noel gật gù, vẫn còn chìm trong dư vị của bữa ăn: ``Ừ, lần sau mà có cá ngừ nữa, em chỉ mong Kuon-senpai là người nấu thôi.''

Cậu con trai từ trong bếp vọng ra: ``Tốt nhất là từ sau đừng ai nhặt con gì lạ về nữa. Anh không muốn văn phòng này biến thành vườn bách thú đâu!''

Bữa trưa kết thúc trong tiếng cười đùa náo nhiệt, cả bốn người quay lại công việc thường ngày. Trong khi Noel đang kiểm tra giấy tờ, Choco chuẩn bị lịch trình mới cho chương trình tiếp theo, Kuon bận rộn với màn hình máy tính, thì Haachama lại nhìn chằm chằm vào cửa sổ đã bị phá hỏng. Kuon thầm cảm thấy một điềm báo chẳng lành: ``\textit{Lần tới\dots\ có khi mình thử nhặt về con cá heo xem sao!}''

\end{document}